\begin{frame}{세 단계로 읽기: 편이는 구조다, 순위는 ``대체로''만}
  \textbf{(1) 정책마다 다른 양으로 벌어진다}
  \begin{itemize}
    \item 고정된 오프셋이면 차를 빼면 그만. 그런데 실물 $-$ 시뮬레이션 차는
      \textbf{정책마다} 달라짐:
      \vskip0.3em
      \small
      $(k_p\!=\!1, k_d\!=\!0)$: $-67.0$ \qquad $(k_p\!=\!4, k_d\!=\!4)$:
      $+103.9$ \qquad $(k_p\!=\!8, k_d\!=\!2)$: $-40.4$
      \normalsize
    \item 전체 평균은 $+5.0$ 에 불과하나 정책별 차이는 약 \textbf{170pt}
      차이 --- ``운''이 아니라 ``마찰 모델''이라는 \textbf{구조}. 데이터를
      \textbf{늘려도 안 사라짐}.
  \end{itemize}
  \vskip0.4em
  \textbf{(2) 그래도 순위는 대체로 유지}: 시뮬레이션 vs 실물 순위
  상관계수 = \textbf{0.76}. 0.05였다면 시뮬레이션 학습은 전혀 쓸모없었을
  것 --- ``약하게나마'' 실물과 상관하는 것이 sim-to-real의 존재 근거.
  \vskip0.4em
  \textbf{(3) 그러나 ``대체로''가 문제}: $(k_p\!=\!1, k_d\!=\!0)$ 는
  시뮬레이션 $-1355.3$ = 20개 중 \textbf{11위}(상위권) $\to$ 실물
  $-1422.3$ = \textbf{20위(꼴찌)}.
  \vskip0.4em
  \begin{center}
    \kq{시뮬레이션 하나로 정책 하나를 골라 ``완료''라고 선언하는 순간,\\
        그 순간부터 격차가 실수로 바뀐다.}
  \end{center}
\end{frame}
